% A.75 candidate. Notes only, unregistered pending complete independent audits.
\subsection{Every modular two-group geometric closure}

\begin{statement}{A.75}{Open}{Candidate: both cusp assertions for modular two-group closures}
Let \(Y\to E\setminus\{O\}\) be an actual connected finite
\'{e}tale cover whose geometric Galois closure has group
\[
 M_{4n}=\langle r,s\mid r^{2n}=s^2=1,\quad
                        srs^{-1}=r^{n+1}\rangle,
 \qquad n\ge4\text{ a power of two}.                 \tag{A.75.1}
\]
Then both assertions of Task A.5 hold for \(Y\) at every
coefficient prime, every induced cusp packet and every permitted
open derived subgroup. In particular its full fixed geometric
pro-\(\ell\) subgroup is the selected cusp inertia, and every
nontrivial arithmetically normalized procyclic pro-\(\ell\)
subgroup is cuspidal.

The proper Galois closure has genus \(n+1\), \(2n\) cusps,
potential good Jacobian reduction of special \(2\)-rank zero
at two, and cusp differences killed by \(2n\). These properties
pass to every actual intermediate in the precise form used below.
The universal question A.45 remains open.
\end{statement}

\begin{proof}[Candidate proof, pending complete independent audits]
Let \(C\to E\) be the actual smooth proper Galois closure.
Enlarge its number field finitely to define the cover, its deck
group, quotient maps, marked points and relevant endomorphisms.
At a place over two enlarge further to obtain good reduction
of the original elliptic curve and the stable models below.
Every deck automorphism then fixes the coefficient valuation
ring. We use an \emph{unmarked proper stable model} in the
ordinary-component contradiction. Marked models enter only
at the final cusp-path comparison.

\paragraph{Actual inertia, genus and the abelian quotient.}
Write \(k=n+1\). The group consists of the pairs
\((a,b)\), with \(a\bmod2n\), \(b\bmod2\), and multiplication
\[
 (a,b)(c,d)=(a+k^b c,b+d),\qquad k^2\equiv1\pmod{2n}.
                                                        \tag{A.75.2}
\]
A rotation and an outside element generate if and only if
the rotation exponent is odd. Two outside elements
\(r^a s,r^b s\) generate if and only if \(a-b\) is odd:
their quotient is \(r^{a-b}\), and all available rotations
otherwise have even exponent. Two rotations cannot generate.
Direct multiplication gives
\[
 [r^a,r^b s]=r^{-na},\qquad
 [r^a s,r^b s]=r^{-n(a-b)}.                            \tag{A.75.3}
\]
Thus every generating pair has commutator
\(\sigma=r^n\), the central involution. The actual cusp
inertia of the punctured elliptic cover is consequently
\[
 I=\langle\sigma\rangle,\quad |I|=2,\quad
 G/I\simeq C_n\times C_2,
 \quad |S_C|=2n,
 \quad g(C)=n+1.                                      \tag{A.75.4}
\]
The genus follows from \(2g(C)-2=2n(2-1)\).
The actual quotient \(E'=C/I\to E\) is connected and
unramified, hence, after choosing an origin above \(O\),
an elliptic isogeny with kernel \(C_n\times C_2\).
In particular it has potential good supersingular reduction
at two, since it is isogenous to the original \(E\).

Here \(\Phi(G)=\langle r^2\rangle\), and only the inclusion
\(I\subset\Phi(G)\) is needed. Squares and commutators lie
in \(\langle r^2\rangle\), the quotient is elementary
abelian of rank two, and the square of \(r\) generates
that subgroup. Alternatively, \(H\Phi(G)=G\) gives two
independent parity vectors in \(H\); the generating-pair
criterion above then gives an odd rotation and an outside
element in \(H\), so \(H=G\). Thus a proper subgroup
satisfies \(HI\ne G\). We do not identify \(I\) with
\(\Phi(G)\).

\paragraph{Compact type and the two possible special ranks.}
Apply Lehr--Matignon's nilpotent ramified-cover theorem to
the actual map \(C\to E\). The base has a smooth good model,
the sole branch section \(O\) is \'{e}tale over the valuation
ring, the deck group is nilpotent, and generic characteristic
zero ramification is tame. Its Theorem 6.13 gives potential
good reduction of \(J_C\) and a tree as the dual graph
of its proper stable reduction.
\cite[Theorem 6.13]{fghLehrMatignon2005}
\cite[p.662]{fghSaidi2007}
No assertion about slopes of marked exceptional components
is inferred from that theorem.

Actual pullback and norm for the quotient by \(\sigma\)
give
\[
 J_C\sim E'\times A,\qquad \dim A=n,\qquad
                     r^n|_A=-1.                       \tag{A.75.5}
\]
Since \(X^n+1=\Phi_{2n}(X)\), the nonzero anti-invariant
factor \(A\) has a unital action of the field
\(\mathbf Q(\zeta_{2n})\). After good reduction its special
rational \'{e}tale \(2\)-divisible Tate space is a module
over \(\mathbf Q_2(\zeta_{2n})\), a field of degree \(n\).
Its dimension is the special \(2\)-rank of \(A\), and
is at most \(n\). It is therefore zero or \(n\).
The latter alternative says that \(A\) is ordinary.
We now exclude it on the actual proper stable model.

\paragraph{The positive-genus components in the ordinary alternative.}
Suppose \(A\) ordinary. The good special Jacobian is the
product of the Jacobians of the positive-genus components
of the proper stable fibre. Its \(\sigma\)-invariant
factor is the supersingular elliptic curve \(E'\), and
its anti-invariant factor is ordinary. A pair of positive-genus
components exchanged by \(\sigma\) is impossible: its
invariant and anti-invariant factors are isogenous to the
same component Jacobian. Components fixed pointwise by
\(\sigma\) contribute to the invariant factor.

Let \(B\) be a component fixed by \(\sigma\) on which
the induced involution is nontrivial, and put
\(h=g(B/\sigma)\), \(d=g(B)-h\). Pullback and norm
identify its invariant Jacobian with that of \(B/\sigma\)
up to isogeny, so the latter has \(2\)-rank zero. Its
anti-invariant factor is ordinary, giving \(f(B)=d\).
For \(r_B\) fixed points with different exponents \(d_P\),
Deuring--Shafarevich and Riemann--Hurwitz give
\[
 d=r_B-1,\qquad
 d=h-1+\frac12\sum_P d_P,
                         \qquad d_P\ge2.             \tag{A.75.6}
\]
\cite[Corollary 1.8]{fghCrew1984}
\cite[Proposition 1.4]{fghYang2024}
These are formulas on smooth component curves. They imply
\(h=0\) and every \(d_P=2\). If \(B\) has positive
genus then \(d>0\), so \(B\) is ordinary, its entire
Jacobian is anti-invariant, and \(\sigma\) has genus-zero
quotient. The genera of all such components sum to \(n\).
The remaining invariant genus is exactly one.

We recall the elementary fourth-root obstruction at its exact
scope. On an ordinary smooth curve with a faithful finite
\(2\)-group action, subtract twice Deuring--Shafarevich
from Riemann--Hurwitz. The resulting terms
\(2|H|(g(B/H)-f(B/H))\) and
\(d_P-2(e_P-1)\) are nonnegative and sum to zero.
Thus each second lower ramification group is trivial.
The coefficient of the quadratic term of a uniformizer
action embeds \(H_1/H_2\) in the additive residue field;
in particular every wild point stabilizer is elementary
abelian. If an automorphism of order four squared to our
\(\sigma\), its induced involution on \(B/\sigma=\mathbf P^1\)
would have one fixed point. That fibre cannot be unramified:
an order-four permutation of a two-point fibre has square
identity, whereas \(\sigma\) exchanges those points.
The fibre is therefore one ramified point fixed by the
order-four automorphism, contradicting the stabilizer condition.
\cite[\S2.2, equation (3) and Exercise 2.12(c)]{fghBouwWewers}

\paragraph{The only possible orbit and its fixed parent nodes.}
Every rotation other than \(1,\sigma\) has a power of
order four with square \(\sigma\). Also
\[
                    (r^a s)^2=r^{(n+2)a}.             \tag{A.75.7}
\]
As \(\gcd(n+2,2n)=2\), the only outside involutions
are \(s\) and \(\sigma s\); every other outside element
has a power of order four with square \(\sigma\).
For any positive-genus anti-invariant component \(B\),
its stabilizer therefore lies in
\[
              H_0=\{1,\sigma,s,\sigma s\}\simeq C_2^2.
                                                        \tag{A.75.8}
\]
This conclusion requires no faithful action of the whole
stabilizer: an element whose square acts as the nontrivial
\(\sigma\) necessarily acts with order four.

The full orbit has at least \(n\) positive-genus components,
while their total genus is only \(n\). Hence there is
exactly one orbit of \(n\) ordinary elliptic components,
each stabilized by \(H_0\), and no other anti-invariant
positive-genus components. The invariant genus one is thus
carried by a unique elliptic component \(B_0\), on which
\(\sigma\) acts identically. It is supersingular.
The whole group fixes its vertex, because it is the unique
component of this kind.

Root the proper stable tree at \(B_0\). All the ordinary
elliptic vertices are in a single \(G\)-orbit and have
the same positive distance from this root. The stabilizer
\(H_0\) fixes the unique path from the root to each such
vertex. In particular it fixes the parent node on the
corresponding elliptic component \(B\).

\paragraph{The stabilizer action has a kernel outside generic inertia.}
A faithful \(C_2^2\) action on an ordinary elliptic curve
cannot have a common fixed point. Indeed a fixed point
rules out a genus-one quotient by Riemann--Hurwitz, so
the quotient is rational. Deuring--Shafarevich would then
require the total contribution \(\sum_P(e_P-1)\) to
be four. A common fixed point contributes three; each
other ramified orbit contributes either two (stabilizer
of order two) or three (stabilizer of order four). The
necessary equation
\[
                         4=3+2a+3b,\qquad a,b\ge0,
                                                        \tag{A.75.9}
\]
is impossible. Applied at the fixed parent node, this
shows that the action of \(H_0\) on \(B\) has a kernel.
The kernel has order two, since \(\sigma\) acts
nontrivially. Write it as \(\langle k_B\rangle\).
It is outside \(I=\langle\sigma\rangle\), and
\(t_B=\sigma k_B\) is also outside \(I\), while its
induced action on \(B\) is precisely \(\sigma\).

Every generic point stabilizer for the actual cover
\(C\to E\) is trivial or the central inertia group \(I\).
Consequently \(t_B\) has no fixed point on the smooth
proper generic curve, including at the cusps.

\paragraph{No smooth special fixed point of this outside involution.}
Suppose a \(\sigma\)-fixed point on \(B\) were smooth
on the whole proper stable special fibre. It would also
be fixed by \(t_B\). Complete the smooth chart as
\(R[[z]]\). The automorphism fixes \(R\), and
\(t_B(z)-z\) has nonzero reduction because its action
on \(B\) is nontrivial. The reduced series vanishes to
a positive finite order at the fixed point. Weierstrass
preparation writes it as a unit times a positive-degree
distinguished polynomial, with no uniformizer factor.
It therefore defines a nonempty finite flat horizontal
fixed locus in this disc.
\cite[Corollary 2.16]{fghVonk2021}
After a finite extension a generic root is an actual
point of the smooth proper curve, by the smooth formal-disc
identification and properness. This contradicts the absence
of generic fixed points of \(t_B\).

The argument is on the \emph{unmarked proper model}.
It is not enough merely to separate a pair of marked
cusps on an inserted rational component: contracting such
a component would return the forbidden smooth fixed point.
On the ordinary elliptic \(B\), formula (A.75.6) gives
exactly two \(\sigma\)-fixed points. Both must therefore
be nodes of the proper stable model.

\paragraph{The proper stable tree now gives a contradiction.}
One of these two nodes is the parent node toward \(B_0\).
The other is a child node. Its nonempty descendant subtree
contains neither the root \(B_0\) nor any ordinary
elliptic vertex, since all the latter have equal depth.
Every positive-genus component has already been accounted
for. Thus this descendant subtree consists entirely of
rational components and has one attachment to its parent.
If it has \(v\ge1\) vertices, its total valence in the
whole tree is \(2v-1\). Proper unmarked stability requires
at least \(3v\), an impossibility. This excludes the
ordinary alternative and proves that the actual good
special Jacobian of \(C\) has \(2\)-rank zero.

\paragraph{Actual cusp functions and all intermediates.}
The kernel of \(E'\to E\) has exponent \(n\). For
any two of its points \(Q_i,Q_j\), the divisor
\(n(Q_i-Q_j)\) is principal. The actual degree-two map
\(h:C\to E'\) has one point \(P_i\) above each
\(Q_i\), with \(h^*Q_i=2P_i\). Pulling back the
principal divisor gives
\[
                     2n(P_i-P_j)\text{ principal}.
                                                        \tag{A.75.10}
\]
For every proper subgroup \(H<G\), the actual quotient
\(D=C/H\) has \([G:HI]\ge2\) cusps, since
\(I\subset\Phi(G)\). It has positive genus, as it
maps finitely and nonconstantly to \(E\). The actual
Jacobian norm makes \(J_D\) an isogeny quotient of
\(J_C\), hence gives potential good reduction and
special \(2\)-rank zero. Norms of the functions in
(A.75.10) kill every cusp difference of \(D\) by
\(2n\). Over an algebraically closed field pushforward
of a point has coefficient one, independently of ramification.

The torsion-divisor criterion of A.41 therefore supplies
a rational joining path to a distinct companion for every
cusp on the appropriate stable marked model of \(D\).
A.42 proves the full fixed-inertia assertion at every
coefficient prime, and A.59 proves all normalizer assertions
for the two-group closure. The original quotient \(H=G\)
is already covered by A.40. Finite extensions define the
algebraic data and reduction; geometric base-point paths
are compared only up to inner conjugacy. A.41--A.42's
restriction argument retains all original packets and
permitted open derived groups. No lifting of fixed elements
through a quotient is used.

The exact finite group and dimension checks are retained
in \path{notes/fgh_A75_certificate.py}. The component,
local fixed-locus and tree arguments establish the result
for every power of two, independently of that finite range.
\end{proof}
